\documentclass[openany]{ctexbook}
\usepackage{amsmath,amssymb}
\usepackage[hidelinks]{hyperref}
\newcommand{\note}[1]{\footnote{译注：#1}}
\title{代数几何基础第五卷预备笔记\\8.18后线性系讨论：第一批}
\author{据用户上传英译编辑稿译出}
\date{2026年11月15日}
\begin{document}
\maketitle
\chapter*{范围与校订状态}
本批从底本第34页“It remains to explain”起，连续译至第35页“$\xi$ generic in $L$”止，包括两条显示公式及夹注。下一处为第35页“From the [(cunutesque?) Fr] point of view”起的域论段落。此为独立增补，不含此前8.18正文，未合并主合订本。
原稿记号和部分数学断言存在待核问题。正文区分可确定的记号订正与不能确定的原稿主张；本批译校不构成对原稿逆命题的证明，也未恢复8.8缺失条件。
\chapter*{8.18后未编号讨论}
还须把8.18的条件(a)、(b)解释清楚（明确写出）。\note{底本印作“1.18”，但紧接8.18并明确指向该命题的(a)、(b)，此处订为8.18。}这使我们一般地考察如下情形。假定 $X$ 在 $k$ 上几何不可约，并给定一个超平面参数空间中的线性子簇 $L$。\note{底本此处写“a linear subariety $L$ of $P$”，随后又用 $L$ 表示被这些超平面包含的线性子簇，并在后文写作 $L^0$。为区分两个不同空间，正文将参数线性系称为 $L$，其基线性子空间称为 $L^0\subset P$；此举不补定后面关于 $F^0$ 的未明约定。}这对应于研究经过给定点 $x$，或在给定光滑点与 $X$ 相切的超平面截面。因此，$L$ 由那些包含 $P$ 中某个线性子簇 $L^0$ 的超平面组成；在上述两种情形中，$L^0$ 分别是一个点，或 $X$ 在光滑点的切空间的像。我们提出如下问题〔底本夹注：“de ravoi”，读法待核〕：对于 $L$ 的泛点（因而对于 $L$ 的某个非空开子集的所有点），$Y$ 是否几何不可约，且
\[
\dim Y=\dim X-1?
\]
这是 Bertini 定理的一个变式，应当放在第3节中〔底本夹注：“j(devrait figurer) 51”〕，并且完全用同一方法处理〔或者，愿意的话，可以将它归结到那里；底本法文夹注“ou, si on veut, s'y ramène”，法文稿第51页〕。
维数问题可以简单地表述为 $f(X)\not\subset L^0$，也就是说，
\[
X'=f^{-1}(P-L^0)
\]
是 $X$ 的稠密开子集。\note{底本写成“$f(X)\notin L-0$”，下一句立即给出上述逆像稠密条件；这里订正为像不包含于中心 $L^0$，不是把像当成一个点。}
令 $Q$ 为原稿所谓“经过 $L^0$ 的超平面的射影空间”〔注意：如果 $L^0$ 由 $E$ 的向量子空间 $F^0$ 定义，原稿写 $Q=P(F^0)$〕，并考虑典范态射（由底本所记 $F^-\to E$ 导出，参见第II章）：\note{这一句的射影对偶约定尚不清楚：原稿同时把 $Q$ 称为超平面参数空间，又在下面使用 $L\simeq Q^\vee$，二者不能不加解释地混同。$F^-$ 的上标也与前面的 $F^0$ 不同。本批保留这些写法并标待核，不据此猜补向量空间的商、对偶或映射。}
\[
u:P-L^0\longrightarrow Q.
\]
再考虑
\[
g=uf':f^{-1}(P-L^0)=X'\longrightarrow Q.
\]
\note{底本第二条显示公式写 $f^{-1}(P-L)$，但前面的 $X'$ 与 $u$ 的定义域均为中心 $L^0$ 的补集；此处统一为 $L^0$。$f'$ 指 $f$ 在 $X'$ 上的限制。}
于是 $L\simeq Q^\vee$，而 $X'_\xi$（$\xi\in L$）这一族，正是相对于态射 $g$ 的超平面截面族。
另一方面，原稿说可以立即看出：对于每个“足够一般的”$\xi\in L$，“$X'$ 在 $X$ 中稠密”，从而“$X'$ 几何不可约，当且仅当 $X$ 几何不可约”。\note{底本在这句话中未写 $\xi$ 下标，但上句正在讨论截面族 $X'_\xi$。可能本意是相应截面中的稠密性，而不是重复全空间 $X'\subset X$；因原稿未明写，正文引述其现存记号，不擅自补成对所有截面的断言。}
在此基础上，Bertini--Zariski 定理说明，只要
\[
\dim g(X')\ge2,
\]
便得到所需的不可约性结论。

（严格说来，原稿提出可将3.1的逆命题写成如下形式：如果 $X$ 几何不可约，则 $Y$ 几何不可约，当且仅当以下两种情况之一成立：或者 $\dim f(X)=2$；或者 $\dim f(X)=1$，且 $f(X)$ 包含在一条定义于 $k$ 上的直线 $D$ 中，并且 $X\to D$ 的泛纤维几何不可约。）\note{底本确实为“$\dim f(X)=2$”，不是“$\ge2$”；它与紧前一句所用的充分条件存在明显张力。这里引述作者拟议的逆命题并标待核，不能把“当且仅当”当成已核实的定理，也不在缺少完整假设和证明时擅自补成一个修订逆命题。}
原稿接着说，这也使得我们在当前带有 $L$ 的版本中，得到 $Y_\xi$ 几何不可约的充要条件，其中 $\xi$ 是 $L$ 的泛点。\note{这句话依赖上面尚待核的逆命题；忠实译出不表示已验证该充要条件。}
\end{document}
